
\ifthenelse{\boolean{orig}}
{4) \textbf{Demostración de } 2)\kern2pt b)}
{\Subsection*{4) Demostración de 2)\kern2pt b)}}

4) 1. Caso en que $\Phi$ es un gerbe.

Puede encontrarse un morfismo étale sobreyectivo de tipo finito $p
\colon X_1 \to X$, tal que $p^*\Phi$ sea un gerbe trivial. Por
descenso, como en 3)\kern2pt2, se ve que basta demostrar el teorema
para $X_1$, $X_1\times_X X_1$ y
\marginpar{403}
$X_1\times_X X_1\times_X X_1$. Se reduce, pues, al caso en que
$\Phi$ es el gerbe de los torsores bajo un haz de grupos
constructible~$F$, cuyas fibras tienen orden primo con las
características residuales de~$S$.

Según SGA~4 IX~2.14, puede encontrarse una familia finita de
morfismos finitos $p_i \colon Z_i \to X$ y, para cada $i$, un
haz de grupos constante finito $C_i$, de orden primo con las
características residuales de~$S$, de modo que se tenga un
\ifthenelse{\boolean{orig}}
{morfismo}
{monomorfismo}
$$
j\colon F \to \prod_{i} p_{i*}C_i=G \quoi.
$$
Sea $\Phi_i$ el stack de los torsores bajo~$C_i$ y $\Psi$ el stack de
los torsores bajo~$G$. Resulta de 2)\kern2pt a) que, tras restringir
$S$ a un abierto no vacío, puede suponerse que los $(C_i,fp_i)$ son
cohomológicamente propios relativamente a~$S$ en dimensión
$\leqslant 1$ y que los stacks $f_*p_{i*}\Phi_i$ son
$1$-constructibles. Resulta entonces de~\Ref{XIII.1.9} que los
$(p_{i*}C_i,f)$ son cohomológicamente propios relativamente a~$S$
en dimensión $\leqslant 1$; lo mismo vale, por tanto, para~$(G,f)$.
Además, como $p_{i*}\Phi_i$ es equivalente al stack de los torsores
bajo el grupo $p_{i*}C_i$ (SGA~4 VIII~5.8), se ve que $f_*\Psi$ es
$1$-constructible.

Como $\R^1f_*G$ es constructible, puede encontrarse un haz
representable por un $Y$-esquema étale de tipo finito~$T$ y un
epimorfismo
$$
a \colon T \to \R^1f_*G
$$
(SGA~4 IX~2.7); además, puede suponerse que la imagen de la sección
identidad de~$T(T)$ está definida por un torsor~$Q$ sobre~$X\times_Y
T=X_T$, bajo~$G|X_T$. Sean $f_T \colon X_T \to T$ el morfismo
canónico, $F_T=F|X_T$, etc. Según 1)\kern2pt b), tras restringir
$S$ a un abierto no vacío, puede suponerse que $(Q/F_T,f_T)$ es
cohomológicamente propio relativamente a~$S$ en dimensión
$\leqslant 0$ y que $f_{T*}(Q/F_T)$ es constructible. Resulta entonces
de~\Ref{XIII.3.1.2} que
\ifthenelse{\boolean{orig}}
{$f_*\Psi$ es constructible.}
{$f_*\Phi$ es $1$-constructible.}

Demostremos que $(F,f)$ es cohomológicamente propio relativamente a
\marginpar{404}
~$S$ en dimensión $\leqslant 1$. Según \Ref{XIII.1.13}~2), basta
demostrar que, para todo esquema $Y_1$ étale sobre~$Y$ y para todo
\ifthenelse{\boolean{orig}}
{torsor $Q_1$ sobre~$X_1=X\times_Y Y_1$,}
{torsor $Q_1$ bajo~$G_1$ sobre~$X_1=X\times_Y Y_1$,}
si
$f_1 \colon X_1 \to Y_1$ es el morfismo canónico, entonces
$(Q_1/F_1,f_1)$ es cohomológicamente propio relativamente a~$S$ en
dimensión $\leqslant 0$. Ahora bien, por definición de~$T$, $Q_1$ es,
localmente para la topología étale de~$Y_1$, imagen inversa de~$Q$,
lo que demuestra nuestra reducción.
